\section{Outline/Summary for Evola's Revolt}

\begin{quotex}
Housekeeping note: \emph{Logres} is the pen name of M. Smallwood, and I appreciate the fellowship and camaraderie of all those here, and at Gnosis group; I thank you all, and also our host. What follows is an outline of the foreword to \emph{Revolt against the Modern World}, an executive summary, Evola's own review, and a few footnotes. There will be a discussion Thursday, starting with a few high points and remarks from the Foreword. Evola claims to establish a metaphysics of history, a possibility denied or inverted by all modern thought. His purpose is to aid those who either can, or will be able, to awake from modern history. We will look closely at the text in the foreword, focus on his insights, and discuss. Food for thought: Contrast Evola's work with the Durant's \textit{History of the World}, O. Spengler's \textit{The Decline of the West}, or A. Toynbee's \textit{History of Civilizations}, or even Marxist anti-metaphysic purporting to give the meaning of history.

\begin{verse}
Between the desire\\
And the spasm,\\
Between the potency\\
And the existence,\\
Between the essence\\
And the descent,\\
Falls the Shadow. 
\end{verse}\flright{\textsc{T. S. Eliot}}

\end{quotex}
\paragraph{I. An Outline for the Foreword in Revolt Against the Modern World}
\subparagraph{A. Justification and Need for the volume}
Amateurish ``intellectuals" have wrung their hands over the ``West" while lacking true principles, retain (due to passions) false principles they ostensibly reject, and get caught up generally in ``contingent and terminal" forms which are the spasms of false principles to begin with. None of these ``reactions" have positive value, only possessing value as a symptom merely, like the sleepwalker whose deep sleep prevents him from noticing he is ambling about on a dreadful height or precipice, subject to the ever-increasing likeliness of tumbling lifelessly down into ruin. This is a protective mechanism which creates an illusory ``limit", allowing the disease to run a terminal course before detection, thereby creating a better climate for the inevitable death of the patient. This process has only lately been detected, and ineffectually, although ancient man was well aware of its possibility and outcome. This demonstrates that the process robbed man of even an awareness of ``true normalcy and health". On the contrary, we must dig ever deeper to uncover the true root of the disease (as opposed to symptoms, since treating symptoms will not lead to a good outcome). A foundation, a center, a principle gives rise to true action, rather than reaction, and men today are incapable of anything but reflexive protest. Even experience teaches that nothing salutary happens in this way: we cannot toss and turn on a bed of agony – the point is awaken and arise. The real starting point is not to protest ``technocracy" or ``consumerism", but to find the root. The real still exists, invisible and mute – we have to acquire eyes and ears to see it. This would be an absolute reference point. This would be a point of understanding modern deviation and simultaneously, mounting an effective defense\footnote{The energies that have been liberated, or are in the course of liberation, are not such as can be reconfined within the structure's of yesterday's world. The very fact that attempts at reaction have referred to those structures alone, which are void of any superior legitimacy, has made the subversive forces all the more vigorous and aggressive. \textit{Evola}.}. ``The only thing that matters is the silent endurance of a few". This will induce liberating crisis in those around them.

\subparagraph{B. Purpose}
This is my task, within my limitations, to contribute to this work of establishing the total decadence of the modern world. It is upon the ruins of real civilization that the ``Modern" has been constructed, and this realization will provide the true foundation for revolt, clarifying what is reacted against and in what name (legitimacy). The cadaverous wisdom of ``Progress" glorifies defeated, decrepit man as the ultimate standard, outside of which is supposedly darkness and nothing. This intransigence and Lie proves the blindness of the Modern.\footnote{``Those whom the gods wish to destroy, they first make mad". Livy.} It is not merely that we gain a certain perspective from realizing the relativity of ``the Modern" in studying what is past, but we also see that from the perspective of what is Past, the Modern simply ``is not": the disappearance of the Modern will itself mean nothing, from the vantage point of Tradition. I call this the ``dualism of civilizations" (as opposed to merely the ``relativity of civilizations"). There is a moral dimension, not merely a historical one, at work here. One possesses Being, the other is built upon illusions. They are different types.\footnote{They are ``unequals". Any argument about which is ``better" is pointless, or incomplete. One is an actual civilization, the other is the opposite of that.}

\subparagraph{C. Thesis}
In the 8th-6th centuries BC, the process of involution and degradation began, when ``History" actually began. There have been several phases following, including the fall of Rome/advent of Christianity, the decay of the feudal order and Empire, and the advent of humanism and Reformation. All of these are Iron Age phenomenon, the ``Wolf Age". Underground forces have become more and more manifest, and the process of decay swifter and swifter, more decisive, more universal. It will seal the collective fate of millions. This is the historical and relative side. The fact that we cannot trace it further shows how deep the process had already gone, \& how a reaction must step outside of history.\footnote{Thucydides, for instance, had a modern conception of history, and is the father of ``history". Ancient man had no concept of ``Time" as we do.} Any civilization that bases itself in the ``temporal" will eventually fall prey to these same forces. It is ``accidental" that true Tradition lies outside of Time, however, that accident demonstrates where we should begin.\footnote{I take his meaning here to be that, from a modern point of view, the descent into history represents a happy evolutionary accident, or at best, a kind of historical twist, which lead to the randomly meaningless but ultimately beneficial situation we find ourselves in today.}

``These things did not just happen once, but they have always been" – a fog cannot obliterate the heavens. The best of ``modern" scholarship is a display of ignorance at best. Although a relative value can be placed on the results of their research if they are used from another perspective, there is always a tendency of other influences to creep in, the presence of which are subtle. Myth, legend, and epic have a high validity in this line of inquiry\footnote{Wagner's \emph{Parzifal} is interesting from this point of view: a look at ancient myths and ``new" ones (the Grail).}. Our perspective is super-individual and non-human.\footnote{He doesn't speak of a ``superman" like Nietzsche, but of what is ``non-human"} These truths either Are or Are Not: they cannot be debated or even discussed (in the modern way). It is only possible to remember. They are constant and central and equivalent, regardless of how difficult to discern or how fragmentary in application. ``Certainty and transcendence and universal objectivity is innerly established" by them. Nothing can destroy them, nor can they be reached by alternate means.

\subparagraph{D. Method}
For the first part of the volume, I intend to identify by correspondences the traces of what is spiritual, true, certain, and objective (which must be made to speak and to appear) which are homologous across all times and places – this will result in one Truth which is invariant, transcendent, and the origin. Having achieved this, in part two, we will proceed (with this knowledge) using induction, ``a discursive approximation of a spiritual intuition", which will clarify the pathways and etiology of decline, as well as inform us as to the greater and lesser degrees the truth was manifested. (So part one is ``seeing" what is of value in the wreck, part two is understanding the trajectory of the wreck through its relation to what is real). This induction of part two is neither eclectic (as in the sense of picking a certain language best suited to subject or person) nor comparative (just as using parallaxes to find the location of star\footnote{Using parallaxes allows someone in a seemingly relative position without hope of measurement to actually triangulate a relative position vis a vis other distant and also relative points. Erastothenes used the variant positions of shadows on obelisks at Alexandria and in Nubia to calculate the exact circumference of the earth.} is not mere comparison). The result will be ``the same one meaning and same one principle". This will lay the foundations for an eventual positive and effective revolt, against what is passively and semi-consciously accepted today. Those who defend the ``concrete" of the modern world ought to be told, not ``Stop!" but ``Go Ahead!": ``the pit must be filled: there is a need for fertilizer for the new tree that will grow out of your collapse"\footnote{Guido de Giorgio quoted here.}.

\subparagraph{E. Result}
This work will provide guiding principles and essential elements – certainly entire volumes more could be written. These forms and ideas are not ``realities", but that which makes ``realities" possible: their value is independent of that which they establish, which will never be ``perfect". Of course there are imperfections – but this argument would eventually invalidate any reality\footnote{It is certainly true that even traditional civilizations did not utterly do away with sin, injustice, or imperfection (view it how you will).}. Indeed, from this point of view, establishing what is normal and healthy may be no less possible today than yesterday, as even in the past Tradition was not lived up to. We are providing here a ``metaphysics of history", allowing us to trace the genesis of the modern world, and at the same time, embrace what is symbolical, super-historical, and normative. The result will be sufficiently enough (but not exhaustive) for those who are either already awake, or possess the possibility of awakening.

\paragraph{II. Executive Summary of Evola's Volume}
Evola delineates his purpose in the Foreword to \emph{Revolt Against the Modern World}, in a manner to obviate, absolutely, any potential misunderstanding. He clearly explains that the work is a beginning, and a foundational contribution to the preservation of Tradition (rather than a complete edifice, which would require ``volumes"). He is laying a foundation in the sense that he uncovers and defends in an absolute and definitive way that which is, strictly speaking, absolutely necessary to the recovery and manifestation of those principles. Astoundingly, it will be a ``metaphysics" of history, a possibility denied or inverted by all modern thought forms (eg., Spengler, Marx, liberalism, degenerate Christian schools of thought, etc.) Because most Western thought involves and invokes a modicum of history \& theory, it is helpful to explain Tradition both as a counterpoint to modern historicism (Part One: the unreducible remainder of the ``Myths"), and also as a non-historical object that has maintained a decisive presence even in a period which has abandoned it in favor of the temporal continuum, as well as denying it clear and full expression (Part 2).

Therefore, Tradition manifests, in a positive sense in the myths and legends and even some historical periods of the Past (Part 1), \& also as deprivation and endurance (a negative sense) during the period of its so-called eclipse (Part 2), the period in which modernity and anti-Tradition enjoy a titulary and technical dominance that saturates Being today and heavily hampers even the possibility of conceiving Life differently from its current distortions. Because of this, his book includes the ``death" of tradition (in the 6th-8th centuries), \& traces the presence in the periods following this immersion in the heightened historical sense (an immersion which springs to life, for instance, in the histories of Thucydides). However, he also seeks knowledge ``beyond" the historical period (Part One). In this way, he not only traces the rise of Tradition's opposite (Revolution in a pure sense), but forcefully proves the reality of Tradition in a period which did not perfectly or even greatly manifest it (in Evola's mind, there is a distinction between the essence of Tradition which manifests, and the vehicles it uses, as well as the degree or particular form to which it perfects itself: this is why he divides his work into two parts, Part 1 being concerned with the Essence, Part 2 with the particular ``story" of how Tradition manifested).

That this is not a virtuoso parlor trick of the intellect is further shown by the perfect interchangeability and equivalent nature of Tradition throughout the historical continuum in various highly contrasting periods (eg., Greek and Aztec), as also by the fact that this very Tradition radiates an energy which something else has twisted and distorted, making possible the very decay which presumes to deny the primordial source of its remaining power and twisted legitimacy. Evola demonstrates all of this with crystal accuracy, sourcing primary texts, and restates the unwritten thesis in language and arguments suitable particularly for the Western seeker, even the one who is interested in comparing the Western tradition with other traditions. Not only does Evola's thesis possess far more explanatory power than its logical competition (it should be remembered that explanatory power alone is practically the sole criteria of certain higher forms of scientific argument)\footnote{It is certainly true that even traditional civilizations did not utterly do away with sin, injustice, or imperfection (view it how you will).}, it is internally coherent in a way which demonstrates its freedom from truncated and merely ideological solutions to the ``problems" of Life. Another way to put this would be that it is not merely analogically compatible with what we find in history and pre-history, but homologous.

You may find it useful to compare the above with what follows below, which is the

\paragraph{III. Editorial Presentation to the First Edition (Evola's Own Anonymous Review)}
``The seemingly polemical title of this work conceals a powerful historical and metaphysical reconstruction, as the basis for the in-depth understanding the greatest problems of the current age.

The author's fundamental idea is the opposition between two types of civilizations, called ``traditional" and ``modern" respectively: the former, based on the values of pure spirituality, aristocracy, and hierarchy; the latter, rooted in the purely human, secular, and contingent element. With a series of summaries, the fundamental meanings are given that used to dominate in the living, believing, knowing, acting, ruling, and self-transcending of ``traditional man": meanings that an in-depth examination of the most varied texts and evidence are shown to be remarkably identical in all the greatest civilizations of the past. In a second part, the processes are made precise that, in a type of fall, led from the traditional world to the modern world. From the exploration of prehistory, in order to which it throws unsuspected light, the investigation proceeds in an historic outline that extends to the birth of the new Russian-American barbarism. The Aryan civilizations, then the Roman and Ghibelline, in the author's presentation appears as great luminous culminations in this sequence of events over millennia and as symbols of perennial current events.

In a conclusion, the author poses the problem of future times. That last world of the book is not so much a ``decline of the West", but instead a profession of heroic faith, precisely a call for the spiritual revolt after having had the courage of looking into the depths and its most remote roots that world of decadence, against which today those who believe in the highest possibilities of our Revolution, fight.

The book, of which a German translation is in preparation, has a totally particular character, in order to join a series of truly revolutionary historical views to a very serious, scientific documentation presented in an incisive, suggestive, and fully accessible style.

Revolt against the modern world is not a book of sterile polemics, but of serious culture, particularly illuminating for those who want to understand, through a series of quick summaries, the meaning of history and the spirit of the greatest civilizations and institutions of the past as the positive basis to achieve to a truly reconstructive action.

It therefore interests anyone who is not indifferent in the face of the great problems that in the present hour assail the people and from which solutions will depend the destinies of the future world and of Western man in particular."

\flrightit{Posted on 2015-08-22 by Logres }

\begin{center}* * *\end{center}

\begin{footnotesize}\begin{sffamily}



\texttt{Alistair Fraser on 2015-08-24 at 18:27 said: }

Logres , i am pleased to be reading yours and the hosts and others illuminating works , and good tidings from scotland , i am adding my own inputs in various areas


\hfill

\texttt{Mark Citadel on 2015-08-25 at 15:17 said: }

I feel I am definitely going to have to read this one again. My highest recommendation to anyone is to read Guenon's Crisis first, its a gentler introduction to the concepts Evola dives into.


\end{sffamily}\end{footnotesize}
